\documentclass[10pt,a4paper]{article}
\usepackage[T1]{fontenc}
\usepackage[utf8]{inputenc}
\usepackage{lmodern,microtype}
\usepackage[margin=18mm]{geometry}
\usepackage{amsmath,amssymb,booktabs,tabularx,array}
\usepackage{xcolor,hyperref,enumitem}
\definecolor{navy}{HTML}{17365D}
\hypersetup{colorlinks=true,linkcolor=navy,urlcolor=navy}
\setlength{\parindent}{0pt}
\setlength{\parskip}{0.55\baselineskip}
\setlist{nosep,leftmargin=*}
\newcolumntype{Y}{>{\raggedright\arraybackslash}X}
\newcommand{\code}[1]{\texttt{\detokenize{#1}}}
\title{\textbf{Reversible Instance Normalization Experiments}\\Experiment Guideline}
\author{Thesis project}
\date{29 August 2026}
\begin{document}
\maketitle

\section{Role of this document}
This is the current implementation and experiment contract. The rejected
historical ECML/AALTD submission is retained only in the workspace-level parent
directory \code{../../../latex/submissions/ECML_AALTD_2026_reject}. It may provide
historical context, but it is read-only, is not being resubmitted, and does not
define the current protocol.

\section{Forecasting setting}
For user $u$ and query date $t$, the model observes
$X_{u,t}=X_u(t-L,t]$ and forecasts $Y_{u,t}=X_u(t,t+H]$.
Chronological regions own target dates. Test1 evaluates future dates for seen
users; Test2 evaluates future dates for held-out users. All compared methods
share architecture, optimizer budget, batch composition, selected query dates,
and original-space MSE evaluation.

\section{Normalization components}
The location and scale of a lookback may be chosen independently:
\[
\ell_x\in\{0,\operatorname{mean}(x),x_L,\operatorname{median}(x)\},\qquad
s_x\in\{1,\operatorname{std}(x),\operatorname{MAD}(x)\}.
\]
The transformed input is inverted after prediction. The core families are:
\begin{center}
\small
\begin{tabularx}{\linewidth}{@{}lYY@{}}
\toprule
Family & Location/scale & Extra parameters \\
\midrule
\code{none} & none & none \\
\code{standard} & training-global mean/std & none \\
\code{mean_only} & window mean / unit scale & none \\
\code{scale_only} & zero / window std & none \\
\code{instance} & window mean/std & none \\
\code{median_mad} & window median/MAD & none \\
\code{revin} & window mean/std & tied symmetric affine \\
\code{min} & window mean/std & untied global input/output affine \\
\bottomrule
\end{tabularx}
\end{center}
For MIN,
\[
\widetilde x=\gamma\frac{x-\ell_x}{s_x}+\nu,\qquad
\widehat y=s_x\left[\alpha
\frac{f_\theta(\widetilde x)-\nu}{\gamma}+\beta\right]+\ell_x,
\]
with identity initialization. MIN is global, not user-personalized.

\section{Normalized backpropagation}
Each applicable forward transform is paired with data-space MSE and normalized
MSE:
\[
\operatorname{nMSE}(\widehat y,y\mid x)=
\frac1H\sum_{h=1}^H
\left(\frac{\widehat y_h-y_h}{\operatorname{std}(x)+\varepsilon}\right)^2.
\]
Only gradient weighting changes; predictions are denormalized and final tables
report original-space MSE. Pairwise comparisons hold seed, model, optimizer,
batching, and budget fixed.

\section{Model implementation and provenance}
The normalization mechanisms are implemented only in
\code{src/proposal/normalizations.py}. Dataset windows, model composition,
ordinary fitting, manifests, summaries, and plots remain respectively in
\code{data}, \code{model_loading}, \code{training}, \code{pipeline},
\code{results}, and \code{visualization}.

PatchTST is source-adapted from \code{yuqinie98/PatchTST} revision
\code{204c21efe0b39603ad6e2ca640ef5896646ab1a9}; DLinear is source-adapted
from \code{cure-lab/LTSF-Linear} revision
\code{0c113668a3b88c4c4ee586b8c5ec3e539c4de5a6}. Their isolated
\code{external_models} packages use the common \code{lags, dim, horizon}
boundary and are byte-identical to TimeTensors. Normalization remains outside
both borrowed architectures.

\section{Centering and transform ablations}
Two appendix variants use non-affine instance normalization:
\begin{itemize}
\item \code{revin_last}: center on the last observation instead of the mean;
\item \code{revin_arcsinh}: apply a reversible $\operatorname{asinh}$ magnitude
      compression after instance normalization.
\end{itemize}
Each is paired with MSE and nMSE training and compared with the matching
\code{instance} control.

\section{Seed and uncertainty contract}
One seed is the single randomness source for a complete run. It fixes the user
partition, data and window sampling, model initialization, and optimization.
No separate partition or component seed is introduced. Reported cross-seed
dispersion intentionally combines every stochastic source present in the run.

\section{Profiles}
\begin{center}
\small
\begin{tabularx}{\linewidth}{@{}lYYY@{}}
\toprule
Mode & Scope & Models & Budget \\
\midrule
\code{test} & Electricity $504{:}168$, narrow core smoke & PatchTST &
seed 1, 2,000 steps \\
\code{small} & Traffic, Electricity, Solar; five settings & PatchTST &
seeds 1--3, 10,000 steps \\
\code{full} & nine datasets; five settings & PatchTST &
seeds 1--3, 10,000 steps \\
\code{ultra} & full grid & PatchTST and DLinear &
seeds 1--3, 10,000 steps \\
\bottomrule
\end{tabularx}
\end{center}
The shared settings are $168{:}24$, $336{:}48$, and $504{:}168$; RevIN also
uses $336{:}96$ and $336{:}720$. Publication profiles share outputs only when
configuration signatures match.

\section{Execution, artifacts, and tables}
Each DGX Slurm front has a matching \code{_selena.slurm} front that sources the
same family workflow. \code{LOGS_ROOT} and \code{OUTPUTS_ROOT} default to
\path{logs/} and \path{outputs/}; Selena overrides them to
\path{logs_selena/} and \path{outputs_selena/}. Each front owns one family root:
\code{core}, \code{nmse}, \code{exotic}, or \code{min}. The ordered relative
identity below the selected output root is
\begin{verbatim}
family/dataset/L_H/backbone/normalization/loss/run_n/seed_n/
\end{verbatim}
Normalization and loss are separate model configs. Steps, batch size, learning
rate, evaluation cadence, split, and stride are pipeline configs; device and
Slurm placement are runtime-only. Mode selects delayed path subsets and is not
a path or computation-signature field. One seed fixes every stochastic choice.

Each \code{run_n/manifest.json} uses \code{schema_version: 1}, lists declared
parameters, optional plain provenance, required artifacts, and per-seed state,
and has one of the four overall statuses \code{not_run}, \code{running},
\code{interrupted}, or \code{completed}; a seed state may additionally be
\code{ready}. Run identity contains only the schema, identity/model configs,
pipeline/experiment configs, and seeds. It never fingerprints or hashes code, Slurm
files, data, weights, logs, outputs, or directories; those changes are manual
rerun decisions. The schema changes only for a deliberate global contract break,
and \code{new} creates another repeat with unchanged parameters. The default
\code{overwrite_exact} policy skips an identical
completion, resumes an identical interruption, and creates the next run for a
changed pipeline. The overall run remains \code{running} with
\code{ready_at_utc} while finished seed states are \code{ready}; completion is
written immediately after that configuration's producer process returns
successfully. Later configuration or table failures preserve completed runs
and interrupt only unfinished work. Completed manifests
are authoritative, without later file hashing or synchronization checks. Tables
support distinct/latest/average pipeline selection and
selected/latest/distinct/average repeat selection. Explicit pipeline filters
select configurations and must match even with one run;
\code{SELECTED_RUNS.txt} stores only automatic or pinned exact repeats. Reports
live at \code{reports/family/mode} below the selected output root and record requested filters and
obtained input manifests.

Slurm workflows never submit a publisher or execute Git commands. After any
terminal state, including failure, cancellation, or timeout, the user runs
\code{bash publish_job.sh <job-id>} from the project root. The script verifies
\code{main}, sources \path{$HOME/codes/proxy.sh} (or
\code{PROXY_SCRIPT_PATH}), and fast-forward pulls \code{origin/main} before
artifact selection or staging. Lightweight publication selects the exact
local or Selena stdout/stderr pair when a job ID is supplied, all log trees
otherwise, and only aggregate \path{outputs*/reports/} artifacts. Detailed
publication adds all non-binary outputs and diagnostics. Both tiers exclude
\code{*.pt}, \code{*.npy}, and \code{*.cbm}. The script then pushes
\code{origin/main} without opening a pull request. A
non-fast-forward pull stops without a merge commit, and unrelated staged paths
remain outside the publisher commit.

Ninety-four reconstructable configurations were migrated from the earlier
flat-backbone implementation. Their source trees remain under
\path{outputs/archive/legacy_pre_schema_v1_2026-08-07}, but the migrated
\path{outputs/core} manifests also predate the pinned source packages above.
They are retained only as analyzed historical evidence and are not eligible
for current model comparisons. Reproduction must allocate new runs for every
affected PatchTST or DLinear configuration before current tables are used.

Tables report mean and seed standard deviation, explicit row multipliers,
validation-selected policies, and a clearly labeled optimistic test oracle.
The oracle is diagnostic only. Dataset configuration uses shared top-level
fields and project-scoped overrides. For \code{drop_users}, null inherits, an
empty list retains every CSV user, and a nonempty run value replaces the
configured default. After aggregation, \code{missing_values} defaults to
\code{zero}, replacing NaNs with zero; \code{error} rejects them, and
infinite values are always rejected.

\section{Practical validation}
Run external-package forwards, result-format checks, and Slurm workflow tests
locally. Training is a remote cluster task. The first current-package
submission must use a new-run policy so the superseded completed manifests are
not selected as reusable computations.
\end{document}
